Um zu zeigen, dass $\epsilon^{\prime \alpha \beta} = - \epsilon^{\alpha \beta}$, betrachten wir die Bedingung für die Inverse der Lorentztransformation. Die Transformationen sind gegeben durch:

\[
x^{\prime \alpha} = \left(g^{\alpha \beta} + \epsilon^{\alpha \beta}\right) x_\beta
\]

und

\[
x^\alpha = \left(g^{\alpha \beta} + \epsilon^{\prime \alpha \beta}\right) x_\beta^{\prime}.
\]

Setzen wir die erste Gleichung in die zweite ein, erhalten wir:

\[
x^\alpha = \left(g^{\alpha \beta} + \epsilon^{\prime \alpha \beta}\right) \left(g_{\beta \gamma} + \epsilon_{\beta \gamma}\right) x^\gamma.
\]

Da die Transformationen invers zueinander sind, muss gelten:

\[
x^\alpha = \delta^\alpha_\gamma x^\gamma.
\]

Dies führt zu:

\[
\left(g^{\alpha \beta} + \epsilon^{\prime \alpha \beta}\right) \left(g_{\beta \gamma} + \epsilon_{\beta \gamma}\right) = \delta^\alpha_\gamma.
\]

Für infinitesimale $\epsilon^{\alpha \beta}$ und $\epsilon^{\prime \alpha \beta}$ betrachten wir die Terme erster Ordnung:

\[
g^{\alpha \beta} g_{\beta \gamma} + g^{\alpha \beta} \epsilon_{\beta \gamma} + \epsilon^{\prime \alpha \beta} g_{\beta \gamma} = \delta^\alpha_\gamma.
\]

Da $g^{\alpha \beta} g_{\beta \gamma} = \delta^\alpha_\gamma$, folgt:

\[
g^{\alpha \beta} \epsilon_{\beta \gamma} + \epsilon^{\prime \alpha \beta} g_{\beta \gamma} = 0.
\]

Multiplizieren wir diese Gleichung mit $g_{\alpha \delta}$, erhalten wir:

\[
\epsilon_{\delta \gamma} + \epsilon^{\prime}_{\delta \gamma} = 0,
\]

was impliziert, dass $\epsilon^{\prime \alpha \beta} = - \epsilon^{\alpha \beta}$.

Um zu zeigen, dass $\epsilon^{\alpha \beta} = -\epsilon^{\beta \alpha}$, betrachten wir die Erhaltung der Norm:

\[
x^{\prime \alpha} x^{\prime}_\alpha = x^\alpha x_\alpha.
\]

Einsetzen der Transformation ergibt:

\[
\left(g^{\alpha \beta} + \epsilon^{\alpha \beta}\right) x_\beta \left(g_{\alpha \gamma} + \epsilon_{\alpha \gamma}\right) x^\gamma = g^{\alpha \beta} x_\beta x_\alpha.
\]

Die Terme erster Ordnung in $\epsilon$ führen zu:

\[
g^{\alpha \beta} x_\beta \epsilon_{\alpha \gamma} x^\gamma + \epsilon^{\alpha \beta} x_\beta g_{\alpha \gamma} x^\gamma = 0.
\]

Dies vereinfacht sich zu:

\[
\epsilon^{\alpha \beta} x_\beta x_\alpha + \epsilon_{\alpha \beta} x^\alpha x^\beta = 0.
\]

Da dies für beliebige $x^\alpha$ gelten muss, folgt:

\[
\epsilon^{\alpha \beta} = -\epsilon^{\beta \alpha}.
\]

Die Antisymmetrie von $\epsilon^{\alpha \beta}$ ist somit bewiesen.